% generator_I03.tex -- Prop I.3 (Prob): from the greater of two lines, cut off a part equal to the less.
% Figure and proof geometry from jemmybutton byrne-en-latex.tex (CC-BY-SA), wording from 1847 original.
% Build: lualatex generator_I03.tex  ->  generator_I03.pdf
\documentclass[12pt]{article}
\usepackage{geometry}
\geometry{a5paper, margin=1.5cm}
\usepackage[lining]{ebgaramond}
\usepackage{amssymb}
\usepackage{byrne}
\def\mpPre{u := 1cm; textLabels := false;}

\begin{document}

% FIGURE: greater line ABC (AB solid, BC dashed), lesser line EF (blue), circle A through D (blue).
\defineNewPicture{%
    pair A, B, C, D, E, F;%
    path P;%
    numeric r;%
    A := (0, 0);%
    r := 7/4u;%
    B := A shifted (r, 0);%
    C := A shifted (4/3r, 0);%
    D := A shifted dir(30)*r;%
    E := A shifted (7/6r, -1/6r);%
    F := A shifted (7/6r, -7/6r);%
    % Layer-1 assertions: AB = EF = r (the construction goal), B on circle A
    if abs(abs(A-B) - r) > 0.01: errmessage "ASSERT: AB != r"; fi;%
    if abs(abs(A-D) - r) > 0.01: errmessage "ASSERT: AD != r"; fi;%
    byLineDefine(A, B, byblack, SOLID_LINE,  REGULAR_WIDTH);%
    byLineDefine(B, C, byblack, DASHED_LINE, REGULAR_WIDTH);%
    byLineDefine(A, D, byred,   SOLID_LINE,  REGULAR_WIDTH);%
    draw byNamedLineSeq(0)(BC,AB,AD);%
    draw byLine(E, F, byblue, SOLID_LINE, REGULAR_WIDTH);%
    draw byCircle.A(A, D, byblue, 0, 0, 0);%
    draw byLabelsOnPolygon(B, A, D)(OMIT_FIRST_LABEL+OMIT_LAST_LABEL, 1);%
    draw byLabelLineEnd(D, A, 0);%
    draw byLabelLineEnd(C, A, 0);%
    draw byLabelPoint(B, angle(B-A) + 45, 2);%
    draw byLabelsOnPolygon(E, F)(ALL_LABELS, 0);%
}

% STATEMENT
\begin{center}
\textbf{Book~I.\enspace Prop.\ III. Prob.}
\end{center}

\vspace{0.3cm}
\noindent
\begin{minipage}[t]{0.50\linewidth}
\itshape From the greater (\drawUnitLine{AB,BC}) of two given straight
lines, to cut off a part equal to the less (\drawUnitLine{EF}).
\end{minipage}%
\hfill
\begin{minipage}[t]{0.46\linewidth}
\centering
\drawCurrentPicture
\end{minipage}

\vspace{0.5cm}

% PROOF
\begin{center}
Draw $\drawUnitLine{AD} = \drawUnitLine{EF}$~(pr.~I.2);\\[3pt]
describe \offsetPicture{10pt}{0pt}{\drawFromCurrentPicture{%
    draw byNamedLine(AD);%
    draw byNamedCircle(A);%
    draw byLabelLineEnd(D, A, 0);%
    draw byLabelLineEnd(A, D, 0);%
}}~(post.~III),\\[3pt]
then $\drawUnitLine{EF} = \drawUnitLine{AB}$.\\[8pt]
For $\drawUnitLine{AD} = \drawUnitLine{AB}$~(def.~15),\\
and $\drawUnitLine{EF} = \drawUnitLine{AD}$~(const.);\\[3pt]
$\therefore \drawUnitLine{EF} = \drawUnitLine{AB}$~(ax.~I).
\end{center}

\begin{flushright}Q.\ E.\ D.\end{flushright}

\end{document}
